\chapter{Foundational machinery}

\section{The dynamical plane}

The first layer consists of the escape-radius and filled-Julia-set
definitions.  For a quadratic polynomial the critical orbit is the
distinguished orbit, and the finite conditions
\[
 O_N=\left\{c: |c|\leq 2,\quad |p_j(c)|\leq 2\text{ for }j\leq N\right\}
\]
form a decreasing compact outer approximation to $\M$.  This part is already
formalized in
\lean{Mlc.Categorical.Mandelbrot}, including
\[
\bigcap_N O_N=\M.
\]
The corresponding increasing approximation by a uniform orbit bound is also
available.  These set-theoretic identities do not imply local connectivity:
an intersection of compact connected-looking approximants needs additional
component control.

\section{Puzzle pieces and moduli}

The Yoccoz puzzle supplies nested dynamical pieces around the critical point.
The analytic core has the shape
\[
\sum_{n\geq 0}\operatorname{mod}(A_n)=+\infty
\quad\Longrightarrow\quad
\bigcap_n P_n=\{0\},
\]
where $A_n=P_n\setminus P_{n+1}$.  The current Lean theorem
proving the corresponding shrinking statement from the formal modulus
hypothesis is
\begin{center}
\lean{MLC.Quadratic.yoccoz_theorem_conformal}.
\end{center}
Its theorem-facing alias is
\begin{center}
\lean{MLC.Literature.yoccoz_puzzle_shrinkage}.
\end{center}

This is a dynamical-plane result.  To use it for MLC one needs a
phase--parameter theorem that transports the shrinking of dynamical pieces to
shrinking of parameter parapuzzle pieces.  The parameter pieces must be
defined with their ray/equipotential graphs; a frozen translate of a full
Green sublevel is not automatically a Yoccoz parapuzzle.

\section{External rays and conformal coordinates}

Douady--Hubbard's external uniformization, Caratheodory boundary theory,
landing, and the phase--parameter correspondence are separate formal
modules.  The present repository contains partial Böttcher and Green
function developments, but its axiom audit intentionally does not claim a
global external-ray landing theorem.  The review therefore records the
Douady--Hubbard result as an interface until all hypotheses and boundary
extensions are machine checked.

\section{Renormalization}

Renormalization turns a quadratic polynomial into a tower of smaller
polynomial-like maps.  A useful formal tower must carry at least:
\[
\begin{array}{ll}
\text{domains and restrictions} & \text{proper degree-two maps}\\
\text{markings} & \text{critical orbit and puzzle data}\\
\text{moduli} & \text{annulus geometry and lower bounds}\\
\text{combinatorics} & \text{primitive/satellite type and bounds}\\
\text{straightening} & \text{parameter-plane correspondence}.
\end{array}
\]
The files \lean{Mlc.RenormalizationTypes},
\lean{Mlc.MoleculeRenormalizationTower}, and
\lean{Mlc.SatelliteRenormalizationTower} currently provide typed interfaces
for part of this list.  They should be treated as a foundation for the
formalization, not as a completed proof of the corresponding published
theorems.
